\documentclass[pdflatex,sn-mathphys-ay]{sn-jnl}

\usepackage[utf8]{inputenc}%
\usepackage[T1]{fontenc}%
\usepackage{graphicx}%
\usepackage{multirow}%
\usepackage{amsmath,amssymb,amsfonts}%
\usepackage{amsthm}%
\usepackage{mathrsfs}%
\usepackage[title]{appendix}%
\usepackage{xcolor}%
\usepackage{textcomp}%
\usepackage{manyfoot}%
\usepackage{booktabs}%
\usepackage{array}%
\usepackage{longtable}%
\usepackage{calc}%
\usepackage{newunicodechar}%
\providecommand{\tightlist}{\setlength{\itemsep}{0pt}\setlength{\parskip}{0pt}}
\newunicodechar{§}{\S}
\newunicodechar{→}{\ensuremath{\rightarrow}}
\newunicodechar{↔}{\ensuremath{\leftrightarrow}}
\newunicodechar{×}{\ensuremath{\times}}
\newunicodechar{·}{\ensuremath{\cdot}}
\newunicodechar{⊃}{\ensuremath{\supset}}
\newunicodechar{°}{\ensuremath{^{\circ}}}
\newunicodechar{≠}{\ensuremath{\neq}}
\newunicodechar{≥}{\ensuremath{\geq}}
\newunicodechar{≤}{\ensuremath{\leq}}
\newunicodechar{≈}{\ensuremath{\approx}}
\newunicodechar{∼}{\ensuremath{\sim}}
\newunicodechar{±}{\ensuremath{\pm}}
\newunicodechar{∓}{\ensuremath{\mp}}
\newunicodechar{∈}{\ensuremath{\in}}
\newunicodechar{∉}{\ensuremath{\notin}}
\newunicodechar{⊥}{\ensuremath{\perp}}
\newunicodechar{∥}{\ensuremath{\parallel}}
\newunicodechar{∧}{\ensuremath{\wedge}}
\newunicodechar{∨}{\ensuremath{\vee}}
\newunicodechar{¬}{\ensuremath{\neg}}
\newunicodechar{≫}{\ensuremath{\gg}}
\newunicodechar{≪}{\ensuremath{\ll}}
\newunicodechar{⇒}{\ensuremath{\Rightarrow}}
\newunicodechar{⇔}{\ensuremath{\Leftrightarrow}}
\newunicodechar{∞}{\ensuremath{\infty}}
\newunicodechar{←}{\ensuremath{\leftarrow}}
\newunicodechar{∂}{\ensuremath{\partial}}
\newunicodechar{∇}{\ensuremath{\nabla}}
\newunicodechar{√}{\ensuremath{\surd}}
\newunicodechar{∝}{\ensuremath{\propto}}
\newunicodechar{≅}{\ensuremath{\cong}}
\newunicodechar{≡}{\ensuremath{\equiv}}
\newunicodechar{÷}{\ensuremath{\div}}
\newunicodechar{∅}{\ensuremath{\emptyset}}
\newunicodechar{¹}{\textsuperscript{1}}
\newunicodechar{²}{\textsuperscript{2}}
\newunicodechar{³}{\textsuperscript{3}}
\newunicodechar{⁴}{\textsuperscript{4}}
\newunicodechar{⁵}{\textsuperscript{5}}
\newunicodechar{⁶}{\textsuperscript{6}}
\newunicodechar{⁷}{\textsuperscript{7}}
\newunicodechar{⁸}{\textsuperscript{8}}
\newunicodechar{⁹}{\textsuperscript{9}}
\newunicodechar{⁰}{\textsuperscript{0}}
\newunicodechar{«}{``}
\newunicodechar{»}{''}

\raggedbottom

\begin{document}

\title[Supplementary Materials]{Supplementary Materials for "Vitality as the Efficiency of Self-Paid Self-Modeling"}

\author*[1]{\fnm{Alexander} \sur{Andriishin}}\email{alexander@andriishin.com}

\affil*[1]{\orgname{Independent Researcher}, \orgaddress{\city{Moscow}, \country{Russia}}}

\maketitle

\section*{\texorpdfstring{S2.1 Operationalization of \(I_{\text{pred}}\) and \(I_{\text{mem}}\): Full Protocol}{S2.1 Operationalization of I\_\{\textbackslash text\{pred\}\} and I\_\{\textbackslash text\{mem\}\}: Full Protocol}}\label{s2.1-operationalization-of-i_textpred-and-i_textmem-full-protocol}

Both informational quantities in the ratio \(\eta_v = I_{\text{pred}}/I_{\text{mem}}\) --- \(I_{\text{pred}}(t,\tau) = I(M_t; X_E^{[t, t+\tau]})\) (predictivity) and \(I_{\text{mem}}(t) = I(M_t; X_E^{\le t})\) (memory) --- are the mutual information between the internal state \(M_t\) and the environment \(X_E\); they are estimated by one and the same sampling procedure, differing only in the environmental window (the future \([t, t+\tau]\) for \(I_{\text{pred}}\), the past and present \(\le t\) for \(I_{\text{mem}}\)). A direct measurement of either, in the sense of formula (1a) of the main text, requires access to all the internal degrees of freedom of the system and is technically impossible for biological, urban, and cosmic systems. The standardized operationalization fixes four elements for each class of systems.

\begin{enumerate}
\def\labelenumi{(\alph{enumi})}
\item
  \emph{The sensory channel} \(x_t\) --- the observable input states. For the bacterium --- the ligand concentration at the chemotactic receptors; for the city --- the flux of resources through the infrastructure (energy, matter, information); for the biosphere --- the flux of solar radiation and the geochemical flows.
\item
  \emph{The internal channel} \(s_t\) --- measurable summaries of the configuration \(M_t\). For the bacterium --- the methylation level of the adaptation machinery; for the city --- administrative-infrastructural parameters (communication density, decision density, budget profile); for the biosphere --- atmospheric composition and biomass pools.
\item
  \emph{The choice of past/future windows} with \(\tau\) specified. For the bacterium --- the characteristic adaptation time of the signaling loop (methylation), on the order of seconds to minutes --- within the life of a single cell, shorter than the division time (generations refer to the predicted rate of evolutionary adaptation, § S5, and not to the measurement window of \(\eta_v\)); for the city --- a quarter or a year; for the biosphere --- the response time of the biogeochemical cycle (a year to a decade).
\item
  \emph{The method of estimating the mutual information.} A direct computation of \(I(s_t; x_{t+\tau})\) for \(I_{\text{pred}}\) --- and, symmetrically, \(I(s_t; x_{\le t})\) for \(I_{\text{mem}}\) --- requires knowledge of the joint distribution \(p(s, x)\), which is inaccessible for biological, urban, and cosmic systems. Finite-sample estimates are used; the choice of estimator depends on the nature of the data (discrete/continuous, low-/high-dimensional, sample size \(N\)) and is the same for both windows. The three classes of methods used in the present work are:
\end{enumerate}

\emph{Plug-in estimate with bias correction.} The empirical frequencies \(\hat p(s, x) = N(s, x)/N\) are substituted into the formula \(\hat I = \sum_{s, x} \hat p(s, x) \log[\hat p(s, x)/(\hat p(s) \hat p(x))]\). The main problem is the systematic overestimation of MI for small samples (\emph{finite-sample bias}), of order \((m-1)/(2N \ln 2)\) bits, where \(m\) is the number of occupied cells in the histogram. Standard corrections: \textbf{Miller--Madow} (Miller 1955) --- adding to the naïve empirical entropy estimate the term \((m-1)/(2N \ln 2)\) to remove the systematic underestimation, valid for \(m \ll N\); \textbf{Grassberger--Schürmann} (Grassberger 2003) --- an entropy estimate via the digamma function, \(\hat H_{\text{GS}} = \log N - \frac{1}{N} \sum_s n_s \psi(n_s)\), giving smaller bias at intermediate values of \(m/N\). Plug-in is optimal for \textbf{discrete biological data} of low/medium dimensionality (receptor methylation statuses, discrete cell states, \(\dim \le 10\), sample \(N \sim 10^3\)--\(10^5\)): the naturally discrete nature of the signal, the absence of binning problems.

\emph{Kraskov--Stögbauer--Grassberger (KSG; Kraskov et al.~2004).} A \(k\)-nearest-neighbor estimate for \textbf{continuous data}, requiring no binning. The formula: \(\hat I_{\text{KSG}} = \psi(k) + \psi(N) - \langle \psi(n_x + 1) + \psi(n_y + 1) \rangle\), where \(\psi\) is the digamma function and \(n_x\) and \(n_y\) are the number of neighbors within the marginal \(L_\infty\)-neighborhoods whose size is set by the \(k\)-th nearest neighbor in the joint space. It adaptively accounts for the local density of the distribution. The parameter \(k\) controls the \emph{bias--variance trade-off}: smaller \(k\) → less bias, more variance; the authors recommend \(k = 4\). It is applicable to \textbf{continuous time series of medium dimensionality} (\(\dim \le 50\)), sample \(N \sim 10^3\)--\(10^6\) --- the standard for neurophysiological, physiological, and econometric time series.

\emph{Mutual Information Neural Estimation (MINE; Belghazi et al.~2018).} It uses the \textbf{Donsker--Varadhan dual representation} of the KL divergence: \(I(X; Y) = \sup_T \{\mathbb{E}_{p_{XY}}[T(x, y)] - \log \mathbb{E}_{p_X \otimes p_Y}[\exp T(x, y)]\}\), where the supremum is taken over all measurable functions \(T: \mathcal{X} \times \mathcal{Y} \to \mathbb{R}\). The function \(T\) (the ``critic'') is parameterized by a neural network and trained by stochastic gradient descent over mini-batches, on the identity that the MI is bounded below by any \(T\). It is applicable to \textbf{high-dimensional continuous data} (images, LLM embeddings, dimension \(> 10^2\)), where KSG suffers from the curse of dimensionality (the volume of the \(L_\infty\)-neighborhood grows as \(r^d\) and the kNN density estimate degrades). The cost is the absence of guarantees of exact convergence when the critic is undertrained, and substantial computational expense (hours to days on a GPU for a typical sample \(N \sim 10^5\), \(\dim \sim 10^3\)).

\textbf{The recommended default method for each data class is as follows.}

\begin{center}\small
\begin{tabular}{@{} >{\raggedright\arraybackslash}p{0.3333\linewidth} >{\raggedright\arraybackslash}p{0.3333\linewidth} >{\raggedright\arraybackslash}p{0.3333\linewidth}@{}}
\toprule
\begin{minipage}[b]{\linewidth}\raggedright
Data class
\end{minipage} & \begin{minipage}[b]{\linewidth}\raggedright
Method
\end{minipage} & \begin{minipage}[b]{\linewidth}\raggedright
Default parameters
\end{minipage} \\
\midrule
Discrete biological (\(\dim \le 10\)) & Plug-in with Miller--Madow correction & sanity check via KSG with \(k = 4\) \\
Continuous time series (\(\dim \le 50\)) & KSG & \(k = 4\) \\
High-dimensional continuous (\(\dim > 10^2\)) & MINE & critic --- an MLP (multilayer perceptron) with 3 hidden layers of 512 neurons each, training \(\ge 10^4\) batch-updates \\
\bottomrule
\end{tabular}
\end{center}

\textbf{Sensitivity diagnostics.} If two methods disagree by more than 30\% on the same sample --- vary the estimator-specific parameters:

\begin{itemize}
\tightlist
\item
  \emph{KSG:} vary \(k \in \{3, 5, 8, 10\}\) and check for a \emph{plateau} --- a stable value of the estimate over the intermediate range of \(k\). The absence of a plateau (a monotone drift of the estimate with \(k\)) indicates an insufficient sample size or strong local inhomogeneity of the density.
\item
  \emph{Plug-in:} vary the bin size when binning continuous data, and compare the Miller--Madow and Grassberger--Schürmann corrections. A discrepancy between the corrections of \(> 10\%\) is a signal that \(m/N\) is too large and that KSG or MINE is needed.
\item
  \emph{MINE:} vary the critic architecture (number of layers, width), check the \emph{convergence of the bound} over training epochs (a stable value over the last \(\sim 10^3\) batch-updates), and the regularization (gradient clipping against blow-up of the gradients in the exp-term). A non-converging bound, negative estimates, or strong dependence on the random seed is a signal to change the architecture or to switch to KSG for lower-dimensional subprojections.
\end{itemize}

For systems with active sampling (a bacterium with motor feedback onto the gradient, a city with administrative intervention), the operational MI estimate is carried out at a fixed policy \(\pi\): \(I_{\text{pred}}(\pi) = I(s_t; x_{t+\tau} \mid \pi)\). Conditioning on \(\pi\) is needed in order to define the \textbf{per-policy} quantity \(I_{\text{pred}}(\pi)\) --- without it the MI estimate would conflate predictive information with the effect of active control. But conditioning on \(\pi\) \textbf{does not restore} the DPI bound: under a deterministic action \(a = \pi(M_t)\) the internal state \(M_t\) remains causally coupled to the future of the environment, and \(I(s_t; x_{t+\tau})\) may exceed \(I(s_t; x_{\le t})\). The correct bound in the active case is given by the \textbf{directed} (transfer) information, not by conditioning on the policy --- an open item (see main § 2.1).

The proxies are not the definition of \(I_{\text{pred}}\) --- they are its operationalization. The same discipline operates in standard physical measurement practice: temperature is measured not by directly counting molecular velocities but through a calibrated proxy --- the expansion of mercury. A discrepancy between different proxies for one system is a signal to refine the proxy, not to abandon the concept.

\emph{Why predictive, not stored information.} The precedent for a thermodynamic link between predictive information and dissipation is set by the work of Still, Sivak, Bell, and Crooks on the thermodynamics of prediction (Still et al.~2012). There the term \textbf{nostalgia} is introduced for past information uselessly retained by the system that has no predictive power: \(I_{\text{nostalgia}} := I_{\text{stored}} - I_{\text{pred}} \ge 0\) by the DPI; equality is reached when the internal statistic \(M_t\) is minimally sufficient for prediction (the Information Bottleneck optimum (Tishby et al.~1999)). Still et al.~show that \(I_{\text{nostalgia}}\) sets a lower bound on the total dissipation of the system. In a non-stationary environment the gap \(I_{\text{nostalgia}} > 0\) is typical and is associated with excess dissipation; in a stationary one with a fixed \(M < \infty\) the gap persists but is bounded above by the memory capacity (Still et al.~2012). The present framework inherits Still's precedent and adds a key requirement. Both quantities in the ratio \(\eta_v = I_{\text{pred}}/I_{\text{mem}}\) are informational and pertain to one and the same modeling loop of the system \(X\): both \(I_{\text{pred}}\) and \(I_{\text{mem}}\) are the mutual information between the internal state \(M_t\) and the environment \(X_E\). Standing apart from the construction of the ratio itself is the \textbf{self-payment predicate} \(S\): it requires that the free energy physically maintaining the correlation of \(M_t\) with the environment be dissipated by the system \(X\) itself (the thermodynamic anchor --- the Still bound below), and not be supplied by an external loop. It is \(S\), and not the denominator of \(\eta_v\), that separates the present framework from the classical formulation of the thermodynamics of prediction, where the modeling and dissipating loops may be different: the denominator here is informational (\(I_{\text{mem}}\)), while the belonging of the dissipation to a single system boundary is the content of the predicate \(S\).

\subsection*{\texorpdfstring{S2.1.1 \(\eta_v \le 1\): An Unconditional Consequence of the DPI; The Thermodynamic Anchor --- The Still Bound}{S2.1.1 \textbackslash eta\_v \textbackslash le 1: An Unconditional Consequence of the DPI; The Thermodynamic Anchor --- The Still Bound}}\label{s2.1.1-eta_v-le-1-an-unconditional-consequence-of-the-dpi-the-thermodynamic-anchor-the-still-bound}

The bound \(\eta_v = I_{\text{pred}}/I_{\text{mem}} \in [0,1]\) (for \(I_{\text{mem}} > 0\); the case \(I_{\text{mem}} = 0\) is fixed by the convention \(\eta_v := 0\), main § 2.7) is a direct consequence of the data processing inequality (DPI) for \textbf{passive} self-modeling --- it is a purely informational statement, requiring no step through thermodynamics (no derivation via Landauer's principle or an erasure argument). Thermodynamics enters separately, as a lower bound on the dissipation of retaining non-predictive memory (the Still bound). Below, the two lines are explicitly separated.

\textbf{The unconditional bound \(\eta_v \le 1\) (passive case).} Memory is formed from the past of the environment together with internal noise, \(M_t = f(X_E^{\le t}, \zeta_v)\), where \(\zeta_v\) is internal noise carrying no information about the future of the environment beyond what is already contained in the past: \(\zeta_v \perp X_E^{[t, t+\tau]} \mid X_E^{\le t}\). Then \(M_t\) is a function of the past (and of independent noise), and the triple forms a Markov chain \[M_t \;\to\; X_E^{\le t} \;\to\; X_E^{[t, t+\tau]}.\] By the DPI the mutual information does not increase under processing along the chain: \[I_{\text{pred}} = I\!\left(M_t;\, X_E^{[t, t+\tau]}\right) \;\le\; I\!\left(M_t;\, X_E^{\le t}\right) = I_{\text{mem}},\] whence \(\eta_v = I_{\text{pred}}/I_{\text{mem}} \in [0, 1]\) and the informational nostalgia \(\nu = 1 - \eta_v \in [0, 1]\). This is an \textbf{unconditional} statement: it follows from the Markov property (the structure of the dependence of \(M_t\) on the past) and the DPI, and requires no finiteness of memory, no stationarity, no ergodicity, and no Landauer step. The exact premise required is precisely the conditional independence of the noise \(\zeta_v \perp X_E^{[t, t+\tau]} \mid X_E^{\le t}\), and not the stronger requirement of ``the absence of feedback''; the separation of these two conditions follows below in the active case.

Algorithmically, it is convenient to check the sufficiency of the representation \(M_t\) at the level of the joint distribution. \textbf{Step 1} (Shannon level): by the inequality \(I(X; Y) \le \min(H(X), H(Y))\), where \(H(\cdot)\) is the Shannon entropy of a discrete random variable (\(H(X) = -\sum_x p(x) \log_2 p(x)\), in bits with the logarithm to base 2; Shannon 1948), we have simultaneously \(H(M_t) \ge I(M_t; X_E^{[t, t+\tau]}) = I_{\text{pred}}\) and \(H(M_t) \ge I(M_t; X_E^{\le t}) = I_{\text{mem}}\): the entropy of the internal state majorizes both informational quantities. \textbf{Step 2} (minimal representation): the standard link between average Kolmogorov complexity and Shannon entropy (Cover and Thomas 2006; Li and Vitányi 2008) gives \(\langle K(M_t)\rangle \ge H(M_t) - O(1)\), where \(\langle\cdot\rangle\) is the average over the invariant measure \(\mu\) and \(O(1)\) is an additive constant (the expected Kolmogorov complexity of a typical realization of the memory state is no lower than its entropy). Consequently a physical representation \(M_t\) retaining \(I_{\text{mem}}\) bits of correlation with the past requires no fewer than \(I_{\text{mem}}\) bits of memory; the same lower bound for \(I_{\text{pred}}\) is weaker, since \(I_{\text{pred}} \le I_{\text{mem}}\). This technique establishes that both quantities in the ratio \(\eta_v\) are informational and are realizable as the content of one and the same loop \(M_t\); the derivation of \(\eta_v \le 1\) itself, meanwhile, has already been obtained from the DPI above and does not need the coding technique.

\textbf{Relation to the entropy rate of the source (for reference).} For an ergodic source with a positive entropy rate \(h_\mu > 0\), the rate of predictive updating is bounded above: \(\dot I_{\text{pred}} \le h_\mu\) --- also a consequence of the DPI under \(M_t = f(X_E^{\le t}, \zeta_v)\) and stationarity. Equality is reached in the limit of optimal prediction, when the entire entropy of the source is predictively relevant; an i.i.d. source (independent and identically distributed --- independent, identically distributed samples, the past carrying no information about the future) realizes the opposite limit \(h_\mu > 0\) with \(I_{\text{pred}} = 0\). This estimate links the framework to the classical line of Bialek--Nemenman--Tishby (Bialek et al.~2001) and is not used in the derivation of \(\eta_v \le 1\). The technical bridge is provided by Brudno's theorem (1983): for \(\mu\)-almost all trajectories \[K(X_E^{[t-\tau, t]}) = h_\mu \cdot \tau + o(\tau) = H(X_E^{[t-\tau, t]}) + o(\tau),\] where \(o(\tau)\) is a term with \(o(\tau)/\tau \to 0\) as \(\tau \to \infty\) (asymptotically small compared to the window length); the equality links the algorithmic complexity of the trajectory of the \emph{source} to the rate \(h_\mu\) and serves as a bridge to the Bialek scale, not as a premise of the bound \(\eta_v \le 1\).

\textbf{The thermodynamic anchor --- the Still bound.} The bound \(\eta_v \le 1\) is informational and in itself energy-free. Thermodynamics enters separately, as the lower price of retaining non-predictive memory, and is counted \textbf{stepwise}. The quantity central to the bound is the \textbf{instantaneous} nostalgia at the memory-update step \(t \to t+1\), \[I_{\text{nonpred}}^{\text{inst}} := I(M_t; X_{E,t}) - I(M_t; X_{E,t+1}) = I_{\text{mem}}^{\text{inst}} - I_{\text{pred}}^{\text{inst}},\] where \(I_{\text{mem}}^{\text{inst}} = I(M_t; X_{E,t})\) and \(I_{\text{pred}}^{\text{inst}} = I(M_t; X_{E,t+1})\) are the instantaneous memory and predictivity of a single step; the equality \(I_{\text{nonpred}}^{\text{inst}} = 0\) is reached when the internal statistic \(M_t\) is minimally sufficient for prediction (the Information Bottleneck optimum (Tishby et al. 1999)). Still, Sivak, Bell, and Crooks (2012) prove that retaining this non-predictive correlation obliges the thermostat to receive heat: at the step \(t \to t+1\) the dissipation is bounded below by the \textbf{instantaneous} nostalgia, \[\langle W_{\text{diss}}^{t\to t+1}\rangle \;\ge\; k_B T \, I_{\text{nonpred}}^{\text{inst}}\] (equation (14) of Still as an equality for the ``work'' protocol, passing into inequality (18) upon adding a non-negative relaxation term). Over a window \(\tau\) the stepwise prices \textbf{add up}, they do not multiply; the horizon-level cumulative \(I_{\text{nonpred}} := I_{\text{mem}} - I_{\text{pred}}\) must not be substituted into the bound (the precise horizon↔instantaneous mapping --- main § 2.1, bound (2)). Here \(k_B T\) enters as the \textbf{coefficient of a bound from a theorem} --- a multiplier in the inequality --- and \textbf{not} as an exchange rate of energy for bits. The Still bound gives a thermodynamic floor on the retention of nostalgia; it does \textbf{not} derive \(\eta_v \le 1\) via erasure and is logically independent of the informational bound. Each bit of instantaneous non-predictive memory costs the system no less than \(k_B T \ln 2\) of dissipated heat --- the strict thermodynamic meaning of retaining an obsolescing model.

The two non-negativities rest on different premises. The instantaneous \(I_{\text{nonpred}}^{\text{inst}} \ge 0\) requires the \textbf{Markov property of the environment} (\(M_t \perp X_{E,t+1} \mid X_{E,t}\)): under it \(M_t \to X_{E,t} \to X_{E,t+1}\) forms a Markov chain and \(I(M_t; X_{E,t+1}) \le I(M_t; X_{E,t})\) by the DPI. The cumulative (horizon-level) \(I_{\text{nonpred}} = I_{\text{mem}} - I_{\text{pred}} \ge 0\) is \textbf{unconditional} for the passive case (by the Markov property \(M_t \to X_E^{\le t} \to X_E^{[t, t+\tau]}\) above) and enters \textbf{only} the informational line of the derivation of \(\eta_v \le 1\), not the bound.

\textbf{The active case (an open problem).} When the internal state \(M_t\) acts on the environment (the motor feedback of a freely-swimming \emph{E. coli}, active-inference agents), the Markov property of the chain \(M_t \to X_E^{\le t} \to X_E^{[t, t+\tau]}\) breaks: the action \(a = \pi(M_t)\) makes \(M_t\) a causal cause of the future of the environment, and formally \(I_{\text{pred}}\) may exceed \(I_{\text{mem}}\) --- the agent itself creates a correlation with the future that is not in its memory of the past. Conditioning on a fixed policy \(\pi\) does not restore the Markov property (the action \(a = \pi(M_t)\) makes \(M_t\) dependent on the future of the environment). The correct quantity in the active case is the \textbf{directed} (transfer) information \(T_{M \to X_E}\), which measures how much the system's own state \(M\) causally contributes to the future of the environment beyond its past; the full thermodynamic bound for feedback control is given by the generalized second law \(\langle\Sigma\rangle \ge -\langle\Delta I\rangle\) (Sagawa and Ueda 2010; Horowitz and Esposito 2014; Parrondo et al.~2015). This is \textbf{an open problem} (main § 5).

\textbf{The interventional (do) reduction: the legitimacy of passive measurement.} The openness of the active case does not devalue the operational measurement of \(\eta_v\) on an immobilized preparation, because immobilization realizes a \textbf{do-intervention} in Pearl's sense (Pearl 2009). The exogenous clamping of the action \(\mathrm{do}(a)\) --- a controlled perturbation that fixes \(a\) by an external hand --- breaks the causal arrow \(M_t \to a \to X_E\) and restores the Markov chain \(M_t \to X_E^{\le t} \to X_E^{[t, t+\tau]}\); by the DPI \[\eta_v^{\mathrm{do}} := \frac{I\!\left(M_t;\, X_E^{[t, t+\tau]} \,\middle|\, \mathrm{do}(a)\right)}{I_{\text{mem}}} \in [0, 1].\] \emph{That is:} under an exogenously clamped action the loop is passive again and the ratio correctly lies in \([0, 1]\). The numerator here is interventional, but the denominator remains observational \(I_{\text{mem}}\): clamping \(\mathrm{do}(a)\) in the present and the future does not change the already-formed correlation of \(M_t\) with the past of the environment, so \(I\!\left(M_t;\, X_E^{\le t}\,\middle|\,\mathrm{do}(a)\right) = I_{\text{mem}}\), and the DPI under the intervention closes \(\eta_v^{\mathrm{do}} \le 1\), fixing both the numerator and the denominator in one interventional regime. It is essential that \(\mathrm{do}(a)\) is an \textbf{intervention}, not conditioning on an observed policy: passive conditioning on the action of a freely-swimming agent would not restore the Markov property (the path \(X_E^{\le t} \to M_t \to a \to X_E^{[t, t+\tau]}\) remains open), whereas the exogenous clamp severs the arrow into \(a\) from \(M_t\). For \emph{E. coli} the physical realization of \(\mathrm{do}(a)\) is \textbf{immobilization}: a pinned cell does not swim, its internal state does not drive the ligand concentration, and the immobilized FRET protocol (main § 3) measures precisely \(\eta_v^{\mathrm{do}} \in [0, 1]\), not the active quantity. We emphasize the status: \(\eta_v^{\mathrm{do}}\) is \textbf{not a new theorem}, but a standard do-reduction of the active loop to a passive one (Pearl 2009); it only fixes the precise sense in which the passive \([0, 1]\)-measurement is legitimate and does \textbf{not} assert boundedness of the active \(\eta_v\).

\emph{Sagawa's demon and the defense of the thermodynamic anchor.} The same feedback framework points to an apparent counterexample --- but \textbf{not} to the bound \(\eta_v \le 1\) (which is a pure ratio of MI \(I_{\text{pred}}/I_{\text{mem}}\) and follows from the DPI by construction, immune to any thermodynamic considerations), but to the \textbf{thermodynamic anchor} --- the floor of the Still bound on the retention of nostalgia. Erasing memory that is \emph{correlated} with the environment is permissible at a cost \(\langle W\rangle \ge k_B T(\Delta S - I_{\text{meas}})\), i.e., below \(k_B T \ln 2\) per bit; since the retained predictive memory is correlated with the environment by construction (\(I_{\text{pred}}^{\text{inst}} = I(M_t; X_{E,t+1}) > 0\)), such a channel would seem able to cheapen retention below the floor of the bound. The mechanism, however, requires a controller that \emph{first} measures the correlated reference, acquiring \(I_{\text{meas}}\); by the measurement↔erasure symmetry (Sagawa and Ueda 2009) the measurement itself costs \(\ge k_B T \cdot I_{\text{meas}}\), so that over the full cycle (acquiring the correlation + spending it on cheap erasure) the total dissipation remains \(\ge k_B T \ln 2\) per bit. In the self-payment loop the accounting is internal: there is no external demon with free access to the correlation, and the predictive memory itself is maintained by the same self-paid budget, so the correlation cannot simultaneously be spent on cheapening retention without double-counting (a Maxwell demon inside a closed system). This analysis defends the \textbf{thermodynamic anchor} (the floor of the Still bound is not circumvented by cheap correlated erasure) and does \textbf{not} touch the bound \(\eta_v \le 1\), which holds on the DPI independently of thermodynamics and does not need this defense. The full formal bound for an arbitrary update channel \(M_t \to M_{t+1}\) with action remains an open problem; in the active case \(\eta_v\) may not be defined at all (the numerator is not subject to the DPI), so the claim that internal self-payment accounting ``holds \(\eta_v \le 1\) in the active case'' is \textbf{not} made here (main § 5). Symmetrically with this: the borrowed feedback inequalities (Horowitz--Esposito; Hartich--Barato--Seifert; Sagawa--Ueda) hold the active \textbf{accounting} of work and dissipation well-defined, but are satisfied for \textbf{any} bipartite system (including a thermostat with \(S = 0\)) and therefore by themselves do \textbf{not} discriminate self-payment; the demarcation is carried by the predicate \(S\), not by them, while \(\eta_v\) retreats to the passively-defined \(\eta_v^{\mathrm{do}}\) (main § 2.7).

\textbf{The Bennett regime: a caveat.} The Still bound (like any Landauer estimate) pertains to the \textbf{irreversible} regime of implementation. In the Bennett asymptotics of globally reversible computation (Bennett 1973, 1982, 1989) the per-bit price drops below the Landauer one, and the thermodynamic floor on the retention of nostalgia loses interpretive force: Bennett's time/space trade-off (1989) emulates an irreversible computation reversibly at the cost of a logarithmic growth of the simulation memory. This is a pre-existing physical limitation of Landauer's principle, not specific to the present framework; in the Bennett regime the thermodynamic anchor is redefined via the number of logical rather than erased bits --- an operationally different quantity requiring separate calibration. The informational bound \(\eta_v \le 1\) is thereby preserved unchanged (it does not rely on erasure). The operational signature of the irreversible regime is measurable heat release \(\dot Q > 0\) during the execution of an informational operation; the discovery of a system in the Bennett regime indicates the boundary of the domain of definition of the thermodynamic anchor, not a refutation of the program (main § 5). Hence a boundary remark on vitality (a plausible corollary of the construction of the predicate \(S\), not a derived theorem). We distinguish \textbf{stepwise} from \textbf{global} reversibility. Life is selected \textbf{toward} the stepwise Landauer limit: biomolecular machines are near-reversible per step (on the order of \(20\,k_B T\) per operation --- Bennett 1982), and this stepwise efficiency is a \textbf{mark} of vitality, not its negation (Still 2012). The boundary remark concerns \textbf{only global} thermodynamic reversibility of the whole operation (\(\dot Q = 0\) everywhere), which is unattainable for three interrelated reasons, of which (b) suffices on its own. (a) The load-bearing condition (ii) of the predicate \(S\) is the correction of the model's error; correction (kinetic proofreading --- Hopfield 1974; Ninio 1975) is thermodynamically necessarily dissipative (Bennett 1982; Landauer 1961: a perpetually-working useful machine cannot store its entire history, so irreversibility is ineliminable), so that a strictly reversible loop is \textbf{utterly incapable} of correcting an error and does not satisfy (ii). (b) A finite-memory tracker of a drifting environment (\(h_\mu > 0\)) must erase an obsolete bit \textbf{at every step} (the Landauer rate of erasure), which by itself gives \(\dot Q > 0\) stepwise --- independently of the \textbf{horizon-level} optimality \(\eta_v = 1\) over the window: the latter does not entail zero \emph{instantaneous} nostalgia, since under stochastic drift \(I_{\text{nonpred}}^{\text{inst}} > 0\) persists. (c) Global reversibility requires unbounded simulation memory (the time/space trade-off, Bennett 1989) plus quasi-staticity --- infinitely slowly, incompatible with drift. Therefore a strictly reversible self-modeling loop, however predictive, \textbf{does not satisfy the predicate of vitality} (\(S = 0\) by construction, through the failure of (ii)); both branches of (ii) --- ``the growth of obligatory dissipation'' and ``the decay of \(X\)'' (itself an irreversibly-dissipative process) --- go back to one and the same prohibition and are both empty in the globally-reversible limit. The verdict rests on the qualitative fact of the dissipativity of correction, \textbf{not} on the magnitude of \(T_v\) (which in the Bennett regime is in any case redefined via the number of logical bits --- an operationally different quantity, see above). In this limited sense vitality is plausibly \textbf{constitutively irreversible} --- a thesis whose formal demonstration remains an open problem. This boundary remark subsumes § S5.1: a strictly reversible system is in any case not an ``\(S = 1\)'' counterexample, whereas § S5.1 separately establishes that the Bennett regime lies \textbf{outside the domain of definition} of the \(T_v\)-anchor (re-operationalization via the number of logical bits, falsifiable at \(\eta_v^{\text{log}} > 1\)) --- two non-contradictory statements about the same limit.

\section*{S2.4 Cohort Normalization: Extended Discussion and the Sigmoid Alternative}\label{s2.4-cohort-normalization-extended-discussion-and-the-sigmoid-alternative}

\subsection*{Extended Discussion of Reference Cohorts}\label{extended-discussion-of-reference-cohorts}

The specific reference cohorts for the six paradigm cases of Section 3 are fixed as follows.

\emph{Bacteria} (§ 3.5): the cohort of prokaryotic chemotactic systems, \(n \sim 200\) from BiGG/KEGG models (\emph{E. coli}, \emph{B. subtilis}, \emph{V. cholerae}, and the like).

\emph{Cities} (§ 3.6): 120--500 large cities from the Bettencourt dataset (Bettencourt et al.~2007).

\emph{Biosphere} (§ 3.7): the only case for which percentile normalization is impossible for want of a cohort. In its place, normalization by a theoretical upper bound is used. The biosphere-as-a-whole is normalized through specific exergy flux: net primary production (NPP) \(\sim 10^{14}\) W divided by total living biomass \(\sim 5 \cdot 10^{14}\) kg (Bar-On et al.~2018) → specific \(P_D \sim 0.2\) W/kg as the normalizing scale \(T_v^{\text{ref}}\) over the window \(\tau_{\text{year}}\) (see § 3.7). Percentile normalization by a cohort is inapplicable here for want of a cohort of comparable biospheres; normalization by the specific exergy flux \(P_D\) replaces it.

\emph{Large language model} (§ 3.4): the cohort of the most recent generations of foundation models, \(n \sim 50\) since 2020.

\emph{Stars} (§ 3.1): the main-sequence cohort, \(n \sim 10^4\) from the Hipparcos catalogue.

\emph{Crystal and hurricane} (§§ 3.2--3.3): normalization to zero. The structural zero \(T_v = 0\) or \(C_v = 0\) is set not by position within a cohort but by the constructive failure of the corresponding capacity (no own dissipation paying for retention; no retained predictive information). Formally: \(\tilde{T}_v(X) := 0\) is adopted by convention for \(X\) with \(T_v(X) = 0\), and \(\tilde{T}_v(X) := 1\) for \(X\) with \(T_v(X) \ge \max_{\text{cohort}} T_v\); within the support of the reference cohort --- formula (2a) of the main text.

\subsection*{The Sigmoid Alternative}\label{the-sigmoid-alternative}

An alternative to percentile normalization is the sigmoid normalization \(\tilde{T}_v = T_v / (T_v + T_v^{\text{ref}})\) via an independent choice of the reference value \(T_v^{\text{ref}}\) for each capacity. The sigmoid procedure yields an absolute scale and is applicable to unique systems, but it requires justifying \(T_v^{\text{ref}}, C_v^{\text{ref}}, S_v^{\text{ref}}\) from independent physical considerations. The present work uses the percentile procedure. The sigmoid is left as an alternative for applications requiring an absolute scale or work with unique systems.

\subsection*{\texorpdfstring{Worked Example: \(I_v\) for \emph{E. coli} by Percentile Normalization}{Worked Example: I\_v for E. coli by Percentile Normalization}}\label{worked-example-i_v-for-e.-coli-by-percentile-normalization}

An illustrative step-by-step computation of \(I_v(\textit{E. coli})\) closes the gap between the formal definition of \(I_v\) --- formula (2b) with normalization (2a) of the main text --- and the operational procedure. The computation is given as \textbf{a worked example, not a calibrated measurement}: the numbers are a methodological illustration of the procedure; the value \(I_v \sim 0.5\) by percentile of the cohort of prokaryotic chemotactic systems is a standalone result of the present worked example, not a calibrated empirical claim.

\emph{(a) Proxy for \(T_v\) --- specific exergy flux \(P_D \cdot \eta_{\text{ex}}\) in W/g.} Heat output of \emph{E. coli} in stationary phase \(\sim 1\)--\(3\) pW/cell (Liu et al.~2020), cell dry mass \(\sim 3 \cdot 10^{-13}\) g, \(P_D \sim 3\)--\(10\) W/g (that is, \(\sim 3 \cdot 10^{3}\)--\(10^{4}\) W/kg). With an exergetic efficiency \(\eta_{\text{ex}} \sim 0.4\), the specific exergy flux \(P_D \cdot \eta_{\text{ex}} \sim 1\)--\(4\) W/g --- the proxy for \(T_v(\textit{E. coli})\). The unit is W/g (flux per gram of dry mass): all cohort figures for \emph{E. coli} below are given in W/g, whereas the biospheric reference of this section (\(P_D \sim 0.2\) W/kg) is normalized in W/kg. Physically the scales differ by \(\sim\!10^{4}\), and visually close values (W/g versus W/kg) should not be read as a \(\sim\!10\)-fold discrepancy.

\emph{(b) Distribution of the proxy across the cohort.} The cohort of prokaryotic chemotactic systems --- \(n \sim 200\) from BiGG/KEGG. Heat output ranges from \(\sim 0.1\) pW/cell in slow-metabolism autotrophic chemolithotrophs (\emph{Nitrosomonas}, methanogens) to \(\sim 10\) pW/cell in fast-growing heterotrophs (\emph{V. natriegens}). The specific exergy flux, at comparable cell mass, yields a distribution of \(P_D \cdot \eta_{\text{ex}}\) with median \(\sim 1\) W/g and interquartile range \([0.3,\, 3]\) W/g.

\emph{(c) Rank of }E. coli* in the distribution.* At \(P_D \cdot \eta_{\text{ex}}(\textit{E. coli}) \sim 2\) W/g the rank falls in the mid-to-upper part of the distribution, \(\tilde{T}_v \approx 0.6\).

\emph{(d) Proxies for \(C_v\) and \(S_v\).} \(\tilde{C}_v\) is the normalized excess entropy \(E\) (Crutchfield and Feldman 2003) or statistical complexity \(C_\mu\) (Shalizi and Crutchfield 2001) on the methylation time series (see main § 2.3: LZ complexity belongs to the class of entropy-rate proxies \(h_\mu\), not \(I_{\text{pred}}\); for \(C_v\) as a measure of model quality, proxies of the \(I_{\text{pred}}\) class are needed). For \emph{E. coli}, with its adaptive reactivity to chemical gradients, the proxy is estimated somewhat above the cohort median, \(\tilde{C}_v \approx 0.55\) --- an illustrative point; the formal distribution of excess entropy across the cohort of \(n \sim 200\) requires separate work. \(\tilde{S}_v\) is the Leiden modularity of the network of chemotactic regulation (Tar--CheA--CheY--FliM--motor axially-symmetric assembly); the architecture is hierarchically modular, \(Q \sim 0.5\) by the Leiden algorithm (Traag et al.~2019), which in the cohort distribution yields \(\tilde{S}_v \approx 0.4\).

\(\tilde C_v\) and \(\tilde S_v\) for \emph{E. coli} are given as illustrative points in the cohort of \(n \sim 200\) prokaryotic chemotactic systems (BiGG/KEGG); the full distribution of excess entropy / Leiden \(Q\) across this cohort is an open problem of operationalization, lying outside the present worked example. The numerical illustration is methodological, not a calibrated measurement.

\emph{(e) Geometric mean.} \(I_v = \sqrt[3]{0.6 \cdot 0.55 \cdot 0.4} = \sqrt[3]{0.132} \approx 0.51\) --- the paradigm-case value of \(I_v\) for the present worked example (\(I_v \sim 0.5\) by percentile of the cohort). The sensitivity to the choice of proxy within a single class is estimated by varying each \(\tilde{\cdot}_v\) over a range of \(\pm 0.1\): the resulting \(I_v\) remains within the corridor \([0.4,\, 0.6]\), the width of the corridor being an indicator of the stability of the percentile procedure within the support of the reference cohort.

The epistemic status of each number is illustrative. An analogous step-by-step computation for the remaining five paradigm cases requires fixing a cohort with equal discipline, which lies beyond the scope of the present work; the paradigm-case values of \(I_v\) in § 3 are order-of-magnitude estimates, methodologically consistent with the present worked example.

\section*{S3.4 Counterfactual Ablation of the LLM: Full Component-by-Component Analysis}\label{s3.4-counterfactual-ablation-of-the-llm-full-component-by-component-analysis}

We apply the counterfactual procedure of § 2.2 of the main text to the concrete case of a large language model. The candidates for inclusion in the boundary are examined in sequence.

\emph{Parametric weights (static after training).} Subtraction: the model does not respond; a clean failure of the loop. They belong to the boundary of the generator but not to the working loop at the moment of inference (the weights are not updated by error).

\emph{KV-cache (context memory at runtime).} Subtraction: the model loses the conversation context; over the window \(\tau \sim\) session time, \(\Delta(-dF_X/dt) > 0\). It belongs to the boundary on session timescales.

\emph{Executing server (GPU + OS).} Subtraction: computation is impossible. It belongs to the boundary of the mechanism, but the energy is supplied by the external power grid.

\emph{Owner corporation (which decides on training, deployment, payment).} Subtraction: the model in its current form ceases to be updated (no new versions), but the current version continues to run on the residual infrastructure. An external component relative to the current version of the model; it is included in the boundary when considering the evolution of the model lineage over time.

Counterfactual ablation distinguishes four nested boundaries \(B_1 \subset B_2 \subset B_3 \subset B_4\) (see main § 3.4). The verdict at each boundary is carried by the \textbf{structural self-payment predicate \(S\)}, not by the numerical value of \(\eta_v\): according to the pair \((S, \eta_v)\) of the main text, demarcation is given by the sign, while the magnitude \(\eta_v = I_{\text{pred}}/I_{\text{mem}}\) is a within-class metric, defined only after passing the screening \(S = 1\).

At \(B_1\) (weights only) and \(B_2\) (weights + KV-cache + executing server) on the session window --- \(S = 0\) through the simultaneous failure of \textbf{both} self-payment conditions: (i) retention and updating of the model are paid for by the external power grid, the loop's own dissipation for this work being zero; (ii) there is no physical channel of error return into the working loop --- the output tokens do not deplete the carrier of \(I_{\text{pred}}\). The model is meanwhile predictively strong (large \(I_{\text{pred}}\)), but its payment loop is open, and therefore \(\eta_v\) as a within-class metric is undefined at this boundary --- there is no self-payment subject for which to measure it.

At \(B_3\) (the training pipeline: dataset + GPU fleet + operators of reinforcement-learning-from-human-feedback (RLHF) fine-tuning) on the model-generation window, own dissipation on the GPUs appears, and the error-return loop closes --- but \textbf{onto the corporate operators}: the self-payment subject becomes the pipeline, not the model. At \(B_4\) (the corporation as a legal entity, the full LLM-corp) the self-funding loop is fully closed --- \(S = 1\) --- but the object measured becomes the firm, not the model. At the boundaries \(B_3\)--\(B_4\) the model is predictive (large \(I_{\text{pred}}\)), yet the magnitude of the predictive fraction \(\eta_v = I_{\text{pred}}/I_{\text{mem}}\) of the corresponding loop is \textbf{not measured} and may be low because of the volume of non-predictive memory: its estimation requires MI-estimation of \(I_{\text{pred}}\) and \(I_{\text{mem}}\) of the pipeline or corporation and is left to future work.

Counterfactual ablation thus grants the LLM no vital status at any of the chosen boundaries: where \(S = 1\) (boundaries \(B_3\)--\(B_4\)), the self-payment subject turns out to be no longer the model. It operationalizes the intuition on which the diagnostics of § 3.4 of the main text rests: widening the boundary widens the self-payment subject (model → pipeline → corporation), and does not dissolve the diagnosis but reassigns it.

\section*{S3.5 Worked Example, E. coli: Detailed Sensitivity Analysis}\label{s3.5-worked-example-e.-coli-detailed-sensitivity-analysis}

A full step-by-step calculation of the predictive efficiency \(\eta_v = I_{\text{pred}}/I_{\text{mem}}\) for \emph{E. coli}, referenced in main § 3.5. Both the numerator and the denominator are \textbf{informational} quantities, each estimated from time series: \(I_{\text{pred}} = I(M_t;\, X_E^{[t, t+\tau]})\) --- how much the internal state knows about the environment's future, \(I_{\text{mem}} = I(M_t;\, X_E^{\le t})\) --- how much it knows about the environment's past and current trajectory. By the data processing inequality \(I_{\text{pred}} \le I_{\text{mem}}\), so \(\eta_v \in [0, 1]\) by construction; nostalgia \(\nu = 1 - \eta_v\). The calculation is offered as \textbf{a worked example, not a calibrated measurement}: the numbers illustrate measurability and order of magnitude on a minimal linear model; they do not fix the value of a real \emph{E. coli}.

\subsection*{The Model and Full Calculation Parameters}\label{the-model-and-full-calculation-parameters}

Chemotaxis is modeled as a \textbf{linear adaptive tracker} in the passive (immobilized) regime --- this is the standard Sourjik--Berg / Cluzel FRET protocol, in which the cell is fixed and the ligand is controlled externally; the passive regime bypasses the active (feedback) case. Three elements:

\begin{enumerate}
\def\labelenumi{(\alph{enumi})}
\item
  \emph{Sensory channel} --- the log-concentration of the ligand \(s_t\) as an AR(1) process with correlation time \(\tau_E\) (environmental drift). Controlled by microfluidics.
\item
  \emph{Internal channel (memory)} --- the methylation level \(m_t = \sum_{j \ge 1} w_j\, s_{t-j} + \text{noise}\) as a leaky integral of the past ligand (linearized Tu-type adaptation) with weights \(w_j = (1-\beta_v)\beta_v^{j-1}\), where \(\beta_v\) is the adaptation memory. The internal state is \(M_t = m_t\). The output of the signaling loop (kinase activity) is read out via FRET (the CheY-P/CheZ FRET pair; Sourjik--Berg 2002), while the methylation \(m\) is a latent adaptation variable reconstructed from the dynamics of this activity (Cluzel et al.~2000).
\item
  \emph{Windows} of the past (\(X_E^{\le t}\), length \(L\)) and the future (\(X_E^{[t,t+\tau]}\), horizon \(\tau\)).
\end{enumerate}

Since \((s_t, m_t)\) are jointly Gaussian, both \(I_{\text{mem}}\) and \(I_{\text{pred}}\) are computed \textbf{exactly} via the logdet of the covariance matrices. The key observation regarding measurability: \(I_{\text{mem}}\) depends \textbf{only} on the covariances \(\text{Cov}(m, s_k)\) and \(\text{Cov}(s_i, s_j)\), recoverable from the time series (the methylation \(m\) is a latent variable reconstructed from the FRET readout of kinase activity; the ligand \(s\) --- via microfluidics). That is, \(I_{\text{mem}}\) is estimated from the same accessible observables as the calorimetric estimate of dissipation from heat output; the operational burden of the informational denominator is \textbf{no worse} than the thermodynamic one (a direct answer to the concern that an informational quantity is harder to measure than calorimetry).

\subsection*{Illustrative Numbers}\label{illustrative-numbers}

For characteristic parameters (horizon \(\tau = 5\), past window \(L = 60\), methylation noise \(r = 0.3\)), the calculation yields a \textbf{finite} \(I_{\text{mem}}\) and a predictive efficiency of order \(O(0.1)\). A few points (values in nats):

\begin{center}\small
\begin{tabular}{@{} >{\raggedright\arraybackslash}p{0.1667\linewidth} >{\raggedright\arraybackslash}p{0.1667\linewidth} >{\raggedright\arraybackslash}p{0.1667\linewidth} >{\raggedright\arraybackslash}p{0.1667\linewidth} >{\raggedright\arraybackslash}p{0.1667\linewidth} >{\raggedright\arraybackslash}p{0.1667\linewidth}@{}}
\toprule
\begin{minipage}[b]{\linewidth}\raggedright
\(\tau_E\) (drift)
\end{minipage} & \begin{minipage}[b]{\linewidth}\raggedright
\(\beta_v\) (memory)
\end{minipage} & \begin{minipage}[b]{\linewidth}\raggedright
\(I_{\text{mem}}\)
\end{minipage} & \begin{minipage}[b]{\linewidth}\raggedright
\(I_{\text{pred}}\)
\end{minipage} & \begin{minipage}[b]{\linewidth}\raggedright
\(\eta_v\)
\end{minipage} & \begin{minipage}[b]{\linewidth}\raggedright
\(\nu\) (nostalgia)
\end{minipage} \\
\midrule
10 & 0.5 & 0.69 & 0.32 & 0.47 & 0.53 \\
10 & 0.9 & 0.50 & 0.14 & 0.27 & 0.73 \\
30 & 0.9 & 0.63 & 0.34 & 0.54 & 0.46 \\
3 & 0.9 & 0.30 & 0.02 & 0.07 & 0.93 \\
\bottomrule
\end{tabular}
\end{center}

\(I_{\text{mem}}\) is finite (\(0.03\)--\(0.72\) nats over the run), \(\eta_v \in [0, 1]\) by the DPI, with a characteristic value \(\eta_v \sim O(0.1)\). The predictive efficiency as the informational ratio \(I_{\text{pred}}/I_{\text{mem}}\) is a quantity of order \(0.1\): the methylation memory retains a noticeable fraction of the environmental information that is predictively relevant over the horizon \(\tau\), while the remainder carries nostalgia \(\nu = 1 - \eta_v\).

\subsection*{\texorpdfstring{Nostalgia: Sensitivity in \(\beta_v\) and \(\tau_E\)}{Nostalgia: Sensitivity in \textbackslash beta\_v and \textbackslash tau\_E}}\label{nostalgia-sensitivity-in-beta_v-and-tau_e}

Nostalgia grows (\(\eta_v\) falls, \(\nu\) rises) when the adaptation memory \(\beta_v\) \textbf{outgrows} the environmental drift \(\tau_E\): the system integrates already-stale ligand, and the fraction of pure ballast increases. At fixed \(\tau_E = 5\), varying \(\beta_v\) from \(0.3\) to \(0.99\) gives:

\begin{center}\small
\begin{tabular}{@{}lllllll@{}}
\toprule
\(\beta_v\) & 0.3 & 0.6 & 0.8 & 0.9 & 0.95 & 0.99 \\
\midrule
\(\eta_v\) & 0.29 & 0.25 & 0.20 & 0.14 & 0.09 & 0.035 \\
\(\nu\) & 0.71 & 0.75 & 0.80 & 0.86 & 0.91 & 0.965 \\
\bottomrule
\end{tabular}
\end{center}

That is, \(\eta_v \approx 0.29 \to 0.035\) as \(\beta_v = 0.3 \to 0.99\). This gives nostalgia a concrete, falsifiable reading: faster environmental drift or slower adaptation (larger \(\beta_v\) relative to \(\tau_E\)) → more ballast → lower predictive efficiency. The sensitivity analysis is carried out precisely on these parameters (\(\beta_v\) --- the adaptation memory, \(\tau_E\) --- the environmental drift time, the horizon \(\tau\) and the past window \(L\)). The number of receptors or clusters does not directly affect the ratio \(\eta_v = I_{\text{pred}}/I_{\text{mem}}\): scaling the number of independent sensory channels raises numerator and denominator in concert and cancels in the ratio.

\subsection*{Sensitivity to the Remaining Parameters}\label{sensitivity-to-the-remaining-parameters}

The horizon \(\tau\) is varied: at short horizons \(I_{\text{pred}}\) is small (little about the future), and at long horizons it saturates because of the limit of the adaptation memory; \(\eta_v\) stays within one order of magnitude around \(O(0.1)\).

The past window \(L\), for \(L \gtrsim\) a few \(\tau_E\), barely changes \(I_{\text{mem}}\): the contribution of the distant past is exponentially suppressed both by the weights \(w_j\) and by the decay of the ligand autocorrelation \(\propto \beta_v^{|i-j|}\), \(\propto a^{|i-j|}\) with \(a = e^{-1/\tau_E}\). This is precisely what justifies truncating the cumulative to a finite window.

The methylation noise \(r\) sets the absolute level of \(I_{\text{mem}}\) (larger \(r\) → smaller \(I_{\text{mem}}\)), but weakly affects the ratio \(\eta_v = I_{\text{pred}}/I_{\text{mem}}\), since it enters numerator and denominator in concert.

\subsection*{Bottlenecks of the Calculation}\label{bottlenecks-of-the-calculation}

The bottleneck of the calculation is not the number of receptors, but (i) the exact experimental mapping FRET-readout ↔ internal state \(M_t\) --- which methylation observable to identify with the model's memory --- and (ii) the choice of window for \(X_E^{\le t}\) and the truncation of the cumulative to a finite \(\tau_{\text{mem}}\). Both require calibration against experimental chemotaxis data (Tu et al.~2008; Cluzel et al.~2000; Sourjik--Berg 2002) and constitute a direction for future work. The passive (immobilized) protocol bypasses the active case: freely-swimming \emph{E. coli} is a feedback regime (the open item of main § 2.1, the Sagawa--Ueda 2010 bound). The numbers remain \textbf{illustrative} --- the minimal linear model demonstrates the measurability and order of the predictive efficiency; it does not fix the value of a real \emph{E. coli}; the exact experimental mapping is a separate future-work item. The full reproducible calculation with code is provided as part of the Supplementary Materials (Python sources, saved outputs, parameter documentation).

\section*{S4.4 Counterfactual Ablation of the Thermostat: Full Component-by-Component Analysis}\label{s4.4-counterfactual-ablation-of-the-thermostat-full-component-by-component-analysis}

Apply the counterfactual procedure of § 2.2 to a bimetallic thermostat \(X\) with hysteresis \(\Delta T_{\text{hyst}} = 1\) K, controlling a room with thermal inertia \(\tau_{\text{room}} \sim 30\) min. Candidates for inclusion in the boundary of \(X\) are tested one by one.

\emph{Bimetallic strip (sensor and actuator).} Subtraction: the loop is broken, no regulation. \(\Delta(-dF_X/dt) > 0\) on times \(< \tau_d\). Belongs to the boundary.

\emph{Relay contact set (energy switching).} Subtraction: the same --- the loop is broken. Belongs to the boundary.

\emph{Heating element (the actuating link, powered by the external power grid).} Subtraction: the thermal source disappears; the bimetal and the relay remain. However, \(-dF_X/dt\) of the remaining system tends to zero not through a collapse of the loop, but through depletion of its own free energy (the bimetal slowly cools toward the surroundings). Does not belong to the decision-making loop; an external component.

\emph{The power grid as a current source.} Like the heating element: external.

\emph{The setter (external operator).} Subtraction: the setpoint is fixed, regulation is preserved. External.

The thermostat's boundary under counterfactual ablation is \{bimetal, relay\}. The internal channel that could hold predictive information about the room temperature is the deformation of the bimetal; this information pertains to the environment ``room'', not to the environment ``thermostat + room + power source''. The verdict is carried by the \textbf{structural self-payment predicate \(S\)}, not by the numerical value of \(\eta_v\). In both the control task and the trivial task the thermostat yields \(S = 0\), but through failure of different self-payment conditions.

In the control task ``regulate the room temperature'', condition (i) fails: heating the room is paid for by the external power grid, not by the bimetal's own free energy, so that \(E_{\text{actual}}^{\text{own}} = 0\) relative to this task. The loop of decisions paid for by one's own decay is, for the thermostat, incomplete.

In the trivial task ``hold the bimetal's own shape'', condition (ii) fails: the passive thermo-mechanical coupling forms no error-return loop --- when the bimetal's shape deviates, no own dissipation corrects it; the bimetal simply remains in the new shape as a function of the current temperature. The error of the ``model'' is not returned by physical depletion of the carrier.

Result: \(S = 0\) in both self-consistent interpretations, and \(\eta_v\) as a within-class metric is undefined --- there is no self-payment subject that has passed the screening. The divergence between the interpretations does not indicate an ontological transition --- Bruineberg et al.~(2022) formally show that such a transition cannot be made without covertly importing epistemic assumptions. What it indicates is the conventionality of the boundary choice as a methodological instrument: one and the same physical system admits several self-consistent boundaries, and the choice among them is a methodological, not an ontological, discipline.

\section*{S4.4a Rebutting the Criticisms of Andrews 2021 and Williams 2020}\label{s4.4a-rebutting-the-criticisms-of-andrews-2021-and-williams-2020}

Well-known critiques of the FEP bear on \(\eta_v\) asymmetrically.

Andrews (2021) --- ``The math is not the territory'' --- attacks the FEP for the non-derivability of empirical content from the formalism: without additional hypotheses, \(F\) can be constructed post-hoc for any system. In \(\eta_v\) this critique is rebutted by anchoring the numerator and the denominator to measurable \textbf{informational} quantities --- \(I_{\text{pred}}\) (informational proxies, § 2.3) and \(I_{\text{mem}}\) (MI-estimation, § S2.1); the thermodynamic measurements (exergy) fix a separate side, not entering \(\eta_v\) --- the capacity \(T_v\) and the self-payment predicate \(S\). The freedom of post-hoc construction is removed by the requirement of a concrete operationalization.

Williams (2020) --- ``Predictive coding and thought'' --- attacks predictive coding as a neural mechanism, not the FEP as a principle. This critique is neutral with respect to \(\eta_v\): the program does not assert predictive coding as a mechanism, but uses predictive information as a formal measure, methodologically independent of any particular neural architecture. The distinction between predictive coding (Rao and Ballard 1999; Clark 2013 --- a neuro-architectural hypothesis) and predictive information (Bialek et al.~2001 --- an information-theoretic measure) is essential for a consistent defense of the program against blended criticism.

\section*{S4.7 Organisational Closure: Autopoiesis, M,R-Systems, Constraint Closure}\label{s4.7-organisational-closure-autopoiesis-mr-systems-constraint-closure}

An objection that merits separate treatment: is self-payment not a rebranding of the classical line of organisational closure --- the autopoiesis of Maturana and Varela, the M,R-systems of Rosen, the constraint closure of Mossio--Montévil--Longo, the epistemic cut of Pattee? The direct answer: the four lines fix different sides of closure at the qualitative level; \(\eta_v\) adds to this tradition one thing --- a quantitative thermodynamic discipline through the Still bound (2). This is not an alternative and not a substitute, but a separate analytical register laid over the qualitative conceptual work.

Pattee (Pattee 1982, 2001) introduces the \emph{epistemic cut} as the boundary between the physical-dynamical level and the symbolic level of a system, and \emph{semantic closure} as the closure of the system onto its own processing of its own symbols. Hoffmeyer's biosemiotic program (Hoffmeyer 1996) unfolds this line over the sign processes of the living, fixing semiosis as a constitutive feature of the biological. Our construction of counterfactual ablation of the boundary (§ 2.2 main) is the thermodynamic operationalization of this line: the epistemic cut is drawn not phenomenologically but interventionally, through the discrimination of components by their contribution to the own dissipation of holding the model (bound (2)); the biosemiotic territory of sign processes remains beyond the \(\eta_v\) scale.

Maturana and Varela (Maturana and Varela 1980) define an autopoietic system as a network of processes producing components that recursively realize the same network and define the topological boundary through their own production. Self-payment works in a complementary register: the payment for holding predictive information physically belongs to the same system that holds the information. Maturana and Varela describe organisational closure qualitatively, through the circular structure ``production of components by components''; self-payment supplies the quantitative thermodynamic part of the same discipline through the Still bound (2) --- the thermodynamic price of holding the model.

The categorical distinction of self-payment from autopoiesis: autopoiesis formalizes the \textbf{component} of closure (production of components by components); self-payment formalizes the \textbf{information-thermodynamic identity} of what is held and what is paid for. An autopoietic system can, in theory, hold predictive information under externalized error correction (maintenance of a cell culture in vitro: autopoietic metabolically --- it runs its \textbf{own} metabolism, \(E^{\text{own}} > 0\) --- but the homeostasis of the environment and the return of the model's error are held by the experimenter); the verdict will deny it vital status through failure of the self-payment predicate \(S\) \textbf{by condition (ii), not (i)}: the loop of the model's error-return closes onto the incubator, that is, the subject is reassigned to the composite system ``culture + incubator'' (\(S = 0\) for the culture-in-isolation), and not through the denominator of \(\eta_v\) (the informational \(I_{\text{mem}}\) is indifferent to the source of payment). This is a concrete categorical shift, not merely a quantitative one.

Rosen's M,R-systems (Rosen 1991) are a category-theoretic abstraction of closure with respect to efficient causation: an organism is a system in which every causal role (Aristotle's efficient cause) is closed within the network. Self-payment \(\eta_v\) is its thermodynamic reading: closure is realized through the physical loop model↔dissipation, not through a categorial abstraction. An important substantive distinction: in \emph{Life Itself} Rosen argues that closure-to-efficient-causation is non-computable and non-formalizable in machine-theoretic terms. \(\eta_v\) does not contradict this claim. It does not assert that biology is reducible to computation; it asserts that the thermodynamic price of holding the model is a necessary condition of vitality, measurable formally and not claiming to exhaust the phenomenon of life. The scale fixes one obligatory thermodynamic cross-section, leaving open the question of the non-computable aspects of organisational closure.

The constraint closure of Mossio--Montévil--Longo (Mossio et al.~2016) fixes the closure of biological constraints: each biological constraint is closed through the production of another constraint, and the network of such constraints closes upon itself. Self-payment works at the level of the information-energy loop (predictive information ↔ Landauer dissipation), not at the level of the network of constraints; the constructions are orthogonal, not alternative. Constraint closure describes the topology of the mutual production of constraints; self-payment supplies the thermodynamic price that this topology must pay with its own decay.

\emph{Differential test of the scale with respect to organisational closure.} Consider the pair WT \emph{E. coli} / the minimal cell JCVI-syn3.0. Both systems are autopoietic in the sense of Maturana and Varela, possess M,R-closure in the sense of Rosen, and constraint closure in the sense of Mossio--Montévil--Longo: the qualitative criteria of organisational closure do not distinguish them. The \(\eta_v\) scale gives a differential result on \textbf{two separate axes}. (a) The informational fraction \(\eta_v = I_{\text{pred}}/I_{\text{mem}}\): both the numerator \(I_{\text{pred}}^{\text{syn3.0}}\) (a sharply curtailed sensory repertoire, the absence of two-component chemotactic cascades) and the informational denominator \(I_{\text{mem}}^{\text{syn3.0}}\) are smaller --- both quantities are \textbf{informational}, and their ratio characterizes the predictive fraction of held memory. (b) Separately --- the thermodynamic anchor: the capacity \(T_v\) (metabolic flux, \textasciitilde473 genes against \textasciitilde4300) and the self-payment predicate \(S\) characterize the physical price of holding (the Still bound (2)), \textbf{without} entering the denominator of \(\eta_v\). \(\eta_v\) ranks systems within the organisational-closure class through the informational fraction, while the structural count of capacities \(T_v\)/\(S\) proceeds along a separate axis; the specific estimates are illustrative, by available literature (Hutchison et al.~2016, for the genome; metabolic calorimetry for the specific exergy flux \(T_v\)), requiring independent empirical verification.

Summary of the closure line. The four conceptual predecessors fix different aspects of closure: epistemic-symbolic in Pattee, production-component in Maturana and Varela, efficient-causation in Rosen, constraint-network in Mossio--Montévil--Longo. \(\eta_v\) adds to this tradition a quantitative thermodynamic discipline: the obligatory price of holding non-predictive memory in the system's own degrees of freedom (the Still bound (2)). This is not a rebranding of the four lines, but the construction of a quantitative scale over their qualitative conceptual work. The positional relation is symmetric to the six lines of § 4.6 main: each classical program operationalizes one side of the phenomenon; \(\eta_v\) collects the thermodynamic price into a single measurable quantity, without claiming to replace the corresponding qualitative analytics.

\section*{S4.8 The Origin-of-Life Line and Conceptual Predecessors}\label{s4.8-the-origin-of-life-line-and-conceptual-predecessors}

The organisational closure line (§ S4.7) captures the synchronic aspect of vitality --- how an already-formed living system holds itself together. Developing in parallel is the contemporary \textbf{origin-of-life} (OoL) research program, which captures the diachronic aspect --- under what conditions a closed loop first arises in a geochemical system. In Theory in Biosciences and adjacent journals, this program is represented by five relatively autonomous traditions; each operationalizes one side of the abiotic→biotic transition, and each is a conceptual predecessor of one slice of the \(\eta_v\)-postulate. This subsection articulates that relation explicitly --- not as a claim to close the OoL problem, but as a placement of \(\eta_v\) within an already-existing field.

\emph{The Eigen--Schuster hypercycle.} The hypercycle (Eigen 1971; Eigen and Schuster 1979) is autocatalytic closure at the level of replicators with an information-bearing template carrier. A set of replicators is organized into a catalytic cycle in which each member catalyzes the replication of the next; closure of the error-return loop occurs through cross-catalytic dependencies among the cycle's members. \textbf{Error catastrophe} (Eigen 1971) marks the threshold between the preservation and the loss of information in the replicator loop as the mutation rate grows. This threshold is conceptually homologous to the structural transition \(\eta_v = 0 \to \eta_v > 0\) (§ 2.1 main): in both, the issue is a threshold below which an information-bearing structure loses the ability to hold its own model against a noise channel. Holding \(I_{\text{pred}} > 0\) against noise requires the system to self-pay the thermodynamic floor on the non-predictive ballast (the Still bound, § 2.1 main) --- and the hypercycle's catalytic turnover is precisely this payment. The hypercycle gives a pre-cellular unfolding of the \(\eta_v\) numerator: \(I_{\text{pred}}\) is encoded in the replicator repertoire, \(T_v\) in the hypercycle's catalytic turnover, and closure of the error-return loop through cross-catalytic dependencies. \(\eta_v\) complements the hypercycle with quantitative Landauer discipline: to the structural condition of replicator closure with a template carrier it adds the requirement that the free energy holding this closedness be physically paid for by the cycle's own catalytic turnover rather than subjected to purely external pumping.

\emph{Hordijk--Steel RAF networks.} The reflexively autocatalytic and food-generated (RAF) network formalism (Hordijk and Steel 2004; Hordijk 2019) describes the conditions under which a set of chemical reactions forms a collectively autocatalytic closed network over a given food set. Each reaction in the network is catalyzed by at least one of its own products. The entire network can be assembled from the food set in a finite number of steps. RAF is a structural predecessor of the composite \(C_v\) (informational capacity through the topology of the catalytic network) and \(S_v\) (topological capacity as hierarchical modularity): collective autocatalysis is the structural condition for the appearance of a closed loop. However, RAF is a criterion of \textbf{network closedness}, not a criterion of \textbf{who pays for holding it} in the sense of § 2.1 main. A RAF network may satisfy the closedness conditions and at the same time be paid for by an external chemical flux, not paid for by the system's own dissipative work. \(\eta_v\) complements RAF with the thermodynamic discipline of self-payment: to the structural condition of autocatalytic closedness it adds the requirement that the free energy holding this closedness physically belong to the network itself.

\emph{England's dissipative adaptation.} The term ``dissipative adaptation'' was introduced in Perunov, Marsland, and England (2016) ``\emph{Statistical physics of adaptation}'' (Phys. Rev.~X 6:021036); its statistical-physics foundation is England (2013) ``\emph{Statistical physics of self-replication}'' (J. Chem. Phys. 139:121923), which derived a lower bound on heat production for self-replicating systems (via growth rate, internal entropy, durability). The Perunov--Marsland--England formalism describes dissipative adaptation as a process in which, in systems far from equilibrium subjected to an external driving force, the macrostates that are especially efficient at absorbing and dissipating energy from the environment become statistically fixed. This is a direct thermodynamic predecessor of the \textbf{thermodynamic anchor} of the \(\eta_v\)-program (the Still bound) --- of its \(T_v\)-side (thermodynamic capacity as the system's ability to participate in non-equilibrium exchange). However, dissipative adaptation captures the adaptation of a dissipative structure as such, not a prediction about the environment held by that structure: England's formalism has no \(I_{\text{pred}}\) numerator. \(\eta_v\) adds to dissipative adaptation the requirement that the adapted state hold its own model of the environment --- moving from a pure dissipative structure to a structure that pays for its own modeling.

\emph{Set-inclusion \{P-M-E adapted\} ⊃ \{\(\eta_v\) vital\}.} The Perunov--Marsland--England program covers more broadly than \(\eta_v\)-vitality: any system far from equilibrium that reliably absorbs and dissipates the energy of an external driving field qualifies as ``physical adaptation'' --- independently of whether a self-payment loop is present. Perunov and coauthors are explicit: ``there may be many examples of `well-adapted' structures that did not have parents'' (P-M-E 2016, p.~21) --- physical adaptation as a universal phenomenon, separate from the Darwinian one. The empirical illustration of P-M-E 2016 itself is silver nanorods in a light field (Ito et al.~2013), which statistically ``learn'' to match surface plasmon resonances to the frequency of the pumping field; examples of similar adapted but non-replicator dynamics are also documented in polar actin patterns and active liquid crystals. In \(\eta_v\)-terms these systems are \textbf{adapted but not vital}: they pass the P-M-E ``physical adaptation'' test (reliable matching of absorption-dissipation to external pumping) but fail the counterfactual ablation of § 2.2 --- payment is external (the light field, an ATP influx), there is no self-payment loop. Likewise the thermostat of § 4.4: P-M-E-adapted to grid power, but \(\eta_v^{\text{int}}\) is undefined through the simultaneous failure of \(T_v^{\text{int}}\) and \(C_v^{\text{int}}\). The set-inclusion formalizes the separation: P-M-E adaptation is a \textbf{necessary but not sufficient} condition for \(\eta_v\)-vitality. \(\eta_v\) singles out, within the P-M-E-adapted class, the subset with thermodynamic closedness of the self-payment budget.

\emph{Major transitions in evolution.} The Maynard Smith and Szathmáry (1995) program --- updated in Szathmáry (2015) with the explicit inclusion of information-evolutionary arguments --- captures a sequence of key evolutionary transitions. These are the transitions from molecular replicators to chromosomes and cells, from prokaryotes to eukaryotes, from unicellular to colonial and multicellular forms, from individuals to social groups. Each of them introduces a \textbf{new level of subject of the evolutionary process} with a reassignment of replication and error load. In \(\eta_v\)-terms, a major transition is the event of the appearance of a \textbf{new subject of self-payment}: after the transition, the free energy holding the predictive information about the environment closes onto a new aggregate actor (the cell as an actor over molecules, the multicellular organism as an actor over cells). This correspondence opens a research program: does every major transition in the Maynard Smith--Szathmáry sense qualify as a phase transition in the subject of \(\eta_v\) --- with its own numerator \(I_{\text{pred}}^{\text{new}}\) paid for by dissipation at the new level of organization? This question presents a concrete thermodynamic criterion on the hierarchy of evolutionary levels and remains an open problem for future work.

\emph{Geochemical OoL programs.} The programs of Pross (2012), Smith and Morowitz (2016), and Kauffman (2019) treat the abiotic→biotic transition as the discovery of a closed loop over a geochemical flux. Hydrothermal systems, alkaline vents, and metabolic catalysts form a circuit in which part of the free energy of the external geochemical driving force begins to be held within the circuit's own structure. Pross formulates this as a transition from thermodynamic stability to dynamic kinetic stability. Smith--Morowitz --- as a phase transition in the geochemical network with the appearance of a metabolic core. Kauffman --- as the discovery of the adjacent possible through the spontaneous accumulation of catalytic closures. The Martin--Russell alkaline-vent OoL program (Martin and Russell 2003) captures the geochemical transition to chemoautotrophic prokaryotes as a parallel line relative to the metabolism-first Smith--Morowitz and Pross. Walker (2017b), in a review in \emph{Reports on Progress in Physics}, maps OoL as a problem of physics, closing the context of these programs as a landscape.

In \(\eta_v\)-terms, all three geochemical programs (Pross, Smith--Morowitz, Kauffman) articulate the same phase transition. The moment of origin of life is the moment when a geochemical system transitions from \emph{external payment} (the full free energy passes through the system in transit, without being held) to \emph{self-payment}. Part of the free energy closes within the system's own circuit and physically pays for holding a model of the environment. This formulation has the status of a risk-prediction for a future extension of the methodology beyond the stationary regime; within the present work the operational procedure for measuring \(\eta_v\) on a pre-biotic geochemical substrate is not defined and constitutes an open problem. The counterfactual ablation of the boundary (§ 2.2 main) and the MI-estimation of \(I_{\text{pred}}\) through the sensory channel apply to already-formed biological systems. Their transfer into the pre-biotic register requires the development of proxies for \(I_{\text{pred}}\) and self-payment in an evolving geochemical network.

\emph{The parallel Chirumbolo--Vella 2026 program.} Chirumbolo and Vella (2026) in \emph{Theory in Biosciences} formulate origin of life as an information-dissipative process with water-mediated informational entropy gradients. Their framework is a conceptual neighbor of ours through the shared ontology of informational dissipation; the difference lies in the operational measure: they leave the scale qualitative, whereas \(\eta_v\) sets a quantitative Landauer discipline over the same intuition about self-organizing dissipative systems.

\emph{Summary paragraph.} Five programs --- the Eigen--Schuster hypercycle, RAF networks, dissipative adaptation, major transitions, and geochemical OoL --- each operationalize one side of the phenomenon of the emergence and maintenance of living systems. The corresponding sides: replicator-level autocatalytic closure with a template carrier (hypercycle), structural closedness of the catalytic network (RAF), dissipative adaptation, the evolutionary subject, and the geochemical phase transition. \(\eta_v\) does not compete with any of them and does not claim to close the origin-of-life problem; instead, the scale unifies them into a single dimensionless measure through the self-payment postulate. This is a thermodynamic operationalization of the condition that each of the five programs articulates from its own side: the positivity of \(\eta_v\) is a necessary condition for the moment of origin of life in any of them. Further refinement of each of these correspondences is an open research problem.

\section*{S4.9 Positioning within Contemporary Demarcation Debates}\label{s4.9-positioning-within-contemporary-demarcation-debates}

Beyond the lines of organisational closure (§ S4.7) and origin of life (§ S4.8), the \(\eta_v\)-program explicitly relates to six persistent lines of contemporary philosophy of biology that capture the demarcation difficulty of the concept of ``life.'' Four of them (Cleland, Mariscal--Doolittle, Bich--Green, Pradeu) are outlined in the first paragraph of § 1 main, two (Krakauer et al., Ruiz-Mirazo--Moreno) are added here; below is the positional diagnostic of how each of the six relates to the operational scale \(\eta_v\). The subsection's main methodological thesis: \(\eta_v\) does not offer yet another definition of life, but contributes to the ongoing debate a formal quantitative scale on which the indicated lines can work jointly.

\emph{Cleland's anti-definitionism.} Cleland (2019) argues that the search for a definition of life through lists of features is methodologically barren and has for decades failed to converge. Instead of a dictionary definition, what is needed is a ``theory of life'' --- a theoretical framework in which life occupies a position determined by explanatory structure rather than by list-based demarcation. The empirical background of this line is the Trifonov (2011) registry, which recorded a lexical analysis of 123 definitions of life without convergence to a common core; Cleland's anti-definitionism is a methodological reaction to this material. The \(\eta_v\)-program \emph{joins} this diagnosis: feature-based definitions break down at boundary crossings --- the biosphere, an LLM-corp, a city, and a minimal cell present different combinations of features, not subsumable under a single assortative registry. \(\eta_v\) is not yet another definition, but an operational discriminator: a dimensionless scale on which each system receives a position through the measurable ratio of predictive information to total retained memory (\(\eta_v = I_{\text{pred}}/I_{\text{mem}}\)), without claiming to exhaust the phenomenon of life. The difference. Cleland's anti-definitionism declares the methodological barrenness of a feature-based criterion; \(\eta_v\) offers not a feature-based definition but an operational discriminator with an explicit falsification regime. This is not a theory of life in the full sense, but one of the obligatory dimensions of such a theory.

\emph{Mariscal--Doolittle anti-definitionism.} The ``life and life only'' program of Mariscal and Doolittle (2020) is a radical alternative to definitionism: ``life'' refers to a single monophyletic historical lineage (the terran branch descending from LUCA), not to a list of features, and so defining life through features is futile --- one can only point to the lineage itself. This position is substantively compatible with \(\eta_v\): \(\eta_v\) does not claim a feature-based definition of ``life as such'' either, but gives an operational discriminator of synchronic vitality --- the property of a process of self-payment of predictive information, realized \textbf{multiply} (by the chemical loop of a bacterium, the biogeochemical circuit of the biosphere, a synthetic minimal cell, or a hypothetical non-biological architecture with a closed model↔dissipation loop). The difference: anti-definitionism rejects any feature-based invariant of life; \(\eta_v\) offers one measurable thermodynamic invariant for the synchronic dimension of vitality, without identifying it with ``life.'' The program does not claim to unify the diachronic (population-evolutionary) aspects (§ 2.7 main). This is a risky assertion, falsifiable by the observation of a life-like system across three independent synchronic channels with a stable \(\eta_v = 0\) (§ 5 main).

\emph{Bich--Green operational definitions at frontiers.} Bich and Green (2018) directly discuss how operational criteria should behave at the boundaries of biology --- on viruses, protocells, synthetic constructs. \(\eta_v\) is an operational criterion of exactly this form. For a viral system that does not hold a closed self-payment loop and parasitizes on a host, the virion-in-isolation does not hold an autonomous predictive model of the environment (the numerator \(I_{\text{pred}} \to 0\)), and its holding is not paid for by its own dissipation independently of the host (the predicate \(S = 0\)). The scale gives \(\eta_v = 0\) for the virion-in-isolation (through the numerator) and is consistent with the biological convention, which usually does not assign autonomous vitality to a virus. This point is methodological, not phenomenological: \(\eta_v\) behaves in boundary cases as the Bich--Green program requires of an operational definition, without reproducing the dilemmas of list-based definitions. The difference: on a viral system \(\eta_v = 0\) through the numerator; the virus-in-host changes status --- the program must say \textbf{whose} \(\eta_v\) grows upon infection (reassignment of the subject, § 2.2 main). Bich--Green leave the question of the conventionality of the virus boundary open; \(\eta_v\) fixes it through counterfactual ablation.

\emph{Pradeu's process biology.} Pradeu (2018) argues that the organism is methodologically more productive to think of as a \textbf{process}, not as an object --- biological individuality is set through temporal dynamics, not through a static, object-like identity. \(\eta_v\) operationalizes exactly this: the scale is computed on a window \(\tau_d\) of stability and held across that window (§ 2.6 main); vitality is defined not as a property of the system-as-object, but as a characteristic of the system's dynamics across a temporal window. The window \(\tau_d\), the counterfactual ablation of the boundary, and the stationarity of the triad of invariants (§ 2.3 main) --- all three constructions operationalize Pradeu's processual ontology in a thermodynamic register. The difference: Pradeu's process view excludes an identifiable substrate at all moments of time; \(\eta_v\) requires fixing the boundary of \(X\) through counterfactual ablation on each window \(\tau_d\). This is not a reintroduction of objecthood --- counterfactual ablation does not presuppose the permanence of an object, but only local causal-selective closedness on the window.

\emph{Krakauer et al.~information-theoretic individuality.} The information-theoretic operationalization of individuality by Krakauer et al.~in Theory in Biosciences itself (Krakauer et al.~2020) proposes a decomposition of the system-environment through information-propagation channels with three forms of individuality (organismal/colonial/driven). \(\eta_v\) is complementary to their framework: their apparatus captures the informational side of demarcation without the thermodynamic payment; self-payment adds the Landauer anchor of self-payment (the Still bound), making the three-form typology dependent on \textbf{who pays} for holding the information.

\emph{Ruiz-Mirazo--Moreno basic autonomy.} The basic autonomy program of Ruiz-Mirazo and Moreno (2004) is the line connecting autopoiesis with thermodynamic constraint-logic closest to self-payment. The difference: \(\eta_v\) is positioned as a thermodynamic disciplining over basic autonomy through the explicit identity of the boundaries of the model and the dissipation that pays for it.

\emph{Summary paragraph.} The \(\eta_v\)-program is methodologically compatible with the anti-definitionist line of Cleland, the life-as-lineage program of Mariscal--Doolittle, the operational-frontiers approach of Bich--Green, the process-biology position of Pradeu, the information-theoretic framework of Krakauer et al., and the basic-autonomy program of Ruiz-Mirazo--Moreno --- not as an alternative to them, but as a formal quantitative scale on which these lines can work jointly. This is a contribution to the ongoing demarcation debate, not the proposal of a new definition of ``what life is.'' The program is an operational discriminator that holds a processual ontology, realized multiply through different proxies of the numerator and the denominator, and behaving at boundaries in agreement with biological convention.

\section*{S5 Falsification Criteria: Full Operational Protocols}\label{s5-falsification-criteria-full-operational-protocols}

Main § 5 contains compressed formulations of the nine testable consequences of § 5 (eight substantively independent lines --- see the Note on numbering below), grouped into three blocks. Here we give the full operational protocols with numerical criteria, p-values, sample sizes, preregistration requirements, and conditional caveats. The structure of S5.1--S5.3 parallels Blocks 1--3 of main § 5; S5.4 contains a summary Lakatosian passport of the program.

\emph{Note on numbering.} The nine numbered items of the three blocks reduce to eight substantively independent lines. Item 1 of Block 2 (the statistical aspect --- violation of monotonicity) and item 1 of Block 3 (the constructive aspect --- the power-law dependence, Prediction 1) are two formulations of a single risk-stake about the rate of evolutionary adaptation.

\subsection*{S5.1 Block 1: Direct Refutations of the Scale (Operational Details)}\label{s5.1-block-1-direct-refutations-of-the-scale-operational-details}

\emph{Block 1 item 1 --- \(T_v\) failure (full protocol).} Observation: a vitality-credited (\(S = 1\)) system retaining nonpredictive memory \(I_{\text{nonpred}} = I_{\text{mem}} - I_{\text{pred}} > 0\) over a window \(\tau_d\), whose Still-mandated retention dissipation is demonstrably supplied \textbf{from beyond} its counterfactual-ablation boundary (§ 2.2 main) --- the floor (2) is paid externally, not out of the loop's own budget. Operational test: (i) measurable heat release \(\dot Q > 0\) when an information operation is carried out in an irreversible realization regime (a reversible Bennett realization yields \(\dot Q \to 0\) in the slow-protocol limit); (ii) counterfactual subtraction of the boundary § 2.2 main shows that the heat flux mandated by bound (2) continues to be supplied by an external source after removal of the loop's own degrees of freedom. Fixing \(T\): the temperature of the thermostat into which the system actually dumps its dissipation --- not that of the source-environment and not that of the control setpoint. Such an observation refutes the \textbf{self-payment thesis of the thermodynamic floor} (not the Still bound (2) itself, which is a theorem of stochastic thermodynamics) and nullifies \(T_v\) as a necessary condition. The information bound \(\eta_v \le 1\) is not attacked by this test: it is an unconditional consequence of the DPI (§ 2.1 main) and is not falsified by thermodynamic observations; the subject of this line is self-payment of the thermodynamic floor, not the structure of the relation \(\eta_v\).

\emph{Bennett regime: caveat.} The Bennett regime of reversible computation (Bennett 1973, 1982) does not serve as a refuting scenario: it lies outside the domain of definition and requires a separate operationalization through the number of logical bits. The Bennett regime is a pre-existing physical constraint of Landauer's principle, recognized long before the present program (Bennett 1973); its exclusion is not specific to the \(\eta_v\)-program (any Landauerian measure faces the same domain-of-definition question). A refutation via a Bennett demonstration is credited if the author of the counterexample first operationalizes the ``number of logical bits'' as the canonical numerator and exhibits \(\eta_v^{\text{log}} > 1\).

\emph{Block 1 item 2 --- \(C_v\) failure (full protocol).} Observation: a vitality-credited (\(S = 1\)) system with positive \(C_v^{\text{op}}\), retaining predictive information about a drifting environment over a window \(\tau_d\), for which the bound-(2)-mandated cost of retaining nonpredictive memory (\(I_{\text{nonpred}} = I_{\text{mem}} - I_{\text{pred}} > 0\)) is \textbf{not attributed} to its own budget (under the verdict \(S = 1\) it is paid externally). Operational separation of the budget: counterfactual subtraction of the boundary § 2.2 main separates the loop's own dissipation from the externally supplied one. This refutes the \textbf{consistency of the self-payment predicate \(S\) with \(C_v\) as a necessary condition} (not the Still bound (2) itself, which is a theorem of stochastic thermodynamics).

\emph{Block 1 item 3 --- \(S_v\) failure (full protocol).} Observation: a system with a closed self-payment loop (\(\eta_v > 0\) by counterfactual ablation § 2.2 main) that simultaneously satisfies two conditions: \(Q < Q_{\text{null}} + 2\sigma\) (modularity indistinguishable from a degree-preserving null model --- the configuration model, which preserves the degree distribution) and \(C(k) \not\sim k^{-\beta_h}\) with \(\beta_h \in [0.5,\, 1.5]\) (absence of the hierarchical signature per Ravasz--Barabási 2003). An ensemble of \(\ge 10\) independent biological subjects is required. A refutation is credited at \(p < 0.01\) against the null hypothesis of random topology, at the stated power and with preregistration of the protocol. This nullifies \(S_v\) as a necessary condition.

\subsection*{S5.2 Block 2: Tests of Proxies and Operationalization (Operational Details)}\label{s5.2-block-2-tests-of-proxies-and-operationalization-operational-details}

\emph{Block 2 item 1 --- violation of monotonicity (full protocol).} Observation: a systematic violation of the monotone dependence of the rate of evolutionary learning on \(\eta_v\). If, in a controlled experiment in experimental evolution (bacterial cultures in a chemostat under varying temperature, population density, and carbon source), systems whose \(\eta_v\) differs by an order of magnitude yield rates of adaptation that are statistically indistinguishable over sufficiently long windows, the central ratio loses its numerical meaning. The constructive test coincides with the protocol of S5.3 (Prediction 1).

\emph{Block 2 item 2 --- vitality across three channels at \(\eta_v = 0\) (full protocol).} Observation: a system that indicates vitality along three independent channels --- rate of adaptation, rate of self-repair, and independent biochemical markers of self-production (in the sense of Darwinian evolution; cf.~the NASA definition of life as a ``chemical system capable of Darwinian evolution'' (Joyce 1994) --- here invoked as a conventional external referent, not as a competing theory of vitality; the shift to a substrate-bound register at this step is deliberate). Yet \(\eta_v = 0\) under all operationalizations of \(I_{\text{pred}}\) and \(I_{\text{mem}}\) through standard domain-specific proxies. This is a criterion independent of \(\eta_v\): rate of adaptation, rate of self-repair, and biochemical markers of self-production are operationalized through measurable quantities not reducible to the numerator of \(\eta_v\). A systematic discrepancy would indicate either a defect of the numerator as an operational proxy for predictive information (the methodological level of falsification, see § 2.7 main), or the untenability of the very concept of vitality as the efficiency of self-modeling (the ontological level).

\emph{Block 2 item 3 --- a conventionally living system at \(\eta_v = 0\) (full protocol).} Observation: a biological system that is unquestionably alive by conventional criteria, for which \(\eta_v\) is strictly zero under all reasonable operationalizations. This refutes the adequacy of the scale as an operational discriminator of vitality.

\subsection*{S5.3 Block 3: Progressive Predictions (Operational Details)}\label{s5.3-block-3-progressive-predictions-operational-details}

\emph{Block 3 item 1 --- Prediction 1, power-law dependence of the adaptation rate (full protocol).} The rate of decline of the divergence between the internal model and the environment, normalized to the complexity of the environment, should scale with \(\eta_v\) as a monotonically increasing function. Concrete protocol: a culture of \emph{E. coli} in a chemostat with a controlled carbon flux (glucose, lactose, acetate); \(\eta_v\) is varied through temperature (25--42 °C, five levels) and population density (three levels); adaptation is the decline of the divergence between the encoded methylation level and the observed gradient over 50 generations. Positive prediction: the correlation coefficient (Pearson, between \(\log r_{\text{adapt}}\) and \(\log \eta_v\)) \(\rho \in [0.4,\, 0.7]\) across 45 lines; \(r_{\text{adapt}} = A \cdot \eta_v^{\alpha}\) with \(\alpha \in [0.3,\, 1.0]\). The range is narrow relative to the null hypothesis \(\alpha = 0\) and relative to the hypothesis \(\alpha \ge 2\) (strong nonlinearity); the expected independent checkpoint is data from Lenski's long-term evolution experiment (LTEE), not used in deriving the lower bound. The upper bound \(\alpha = 1\) is direct proportionality of the adaptation rate to the flux \(I_{\text{pred}}\); the lower bound \(\alpha = 1/3\) follows from the structure of (2b) main: \(I_v\) is the cube root of the product of three capacities, and the minimal dependence \(r_{\text{adapt}} \propto \sqrt[3]{\eta_v}\) corresponds to the contribution of one capacity out of three. The scale is calibrated on the worked example of § 3.5 main (predictive fraction \(\eta_v = I_{\text{pred}}/I_{\text{mem}} \sim O(0.1)\)). Refutation: \(\hat{\alpha}\) significantly outside \([0.3,\, 1.0]\) at \(p < 0.01\) (wrong exponent); \(\hat{\rho} < 0.2\) at the stated power (no detectable trend); or \(\hat{\rho} < 0\) (trend in the opposite direction --- a qualitative refutation).

\emph{Block 3 item 2 --- Prediction 2, monotonicity in the matching window (full protocol).} In the regime of approach to stationarity, when \(I_{\text{pred}}\) is not yet saturated and accumulates faster than the integral dissipation, the stationary formulation predicts: under normalization to a unit dissipation rate (\(\dot{E}_{\text{actual}} = \text{const}\) over a pre-stationary transient with a window \(\Delta\tau \ll \tau_d\)), a monotone increase of \(\eta_v(\Delta\tau)\) over the matching interval is expected. Test on the same 45 \emph{E. coli} lines: measure \(\eta_v(\Delta\tau)\) for \(\Delta\tau\) from fractions of a second to seconds (still \(\Delta\tau \ll \tau_d\)) with independent fixation of \(\dot{E}_{\text{actual}}\). A decrease of \(\eta_v\) with \(\Delta\tau\) in this regime is a counter-argument to the stationary framework as such, not a delegation to the companion work.

\emph{Block 3 item 3 --- Prediction 3, structural threshold \(f^*\) (full protocol).} Consider a class of hypothetical autonomous embodied agents with photovoltaic or battery power supply and a physical actuator with its own thermodynamic budget. For this class a structural threshold \(f^*\) of the internalization fraction of the energy budget (the ratio \(E_{\text{actual}}^{\text{own}}/E_{\text{actual}}^{\text{total}}\)) is predicted. Below the threshold, counterfactual ablation § 2.2 main yields \(\eta_v^{\text{int}}\) undefined (structural failure of the self-payment loop); above it, a measurable \(\eta_v^{\text{int}} > 0\). The numerical localization \(f^* \in [0.1,\, 0.5]\) is a research hypothesis of the program; a formal derivation of the bounds via counterfactual ablation is an open problem. Under the two-level formulation of § 3.4 main the structural predicate counterfactual ablation is \textbf{binary} (a self-paid error-return channel either exists or it does not), whereas \(f^*\) is an \textbf{empirical} threshold on a continuous proxy \(f\), above which closure of the loop is statistically stable; \(f^*\) is a regularity of joint internalization, not the structural decision boundary itself.

\emph{Caveat on the domain of applicability.} The prediction does not apply to current cloud-hosted LLMs without state retention, where \(E_{\text{actual}}^{\text{own}} = 0\) constructively (the power grid and infrastructure are external to the forward pass), \(f = 0\) by construction.

\emph{Differential prediction relative to the AI-safety literature.} AI-safety programs (Hubinger et al.~2019; Carlsmith 2022) predict the behavioral emergence of a self-preservation drive at a \emph{capability} threshold --- independent of the energetic architecture; the \(\eta_v\)-program predicts the structural emergence of the self-payment loop at a \emph{thermodynamic-budget} threshold \(f^*\) --- independent of the behavioral profile. Possible separation scenarios: (a) the capability threshold is reached, the self-payment threshold is not (behavioral power-seeking without structural self-payment); (b) the self-payment threshold is reached, capability is not (structural closure of the loop without behavioral drive). Differential observation: a behavioral absence of power-seeking under an observation of \(f \ge f^*\) refutes the coincidence of the two threshold conditions and confirms \(\eta_v\) as an independent predictor.

\emph{Refutation of Prediction 3.} A systematic observation of a system in the indicated class with \(f \in [0,\, 0.1]\) that passes counterfactual ablation with \(\eta_v^{\text{int}} > 0\), or of a system with \(f > 0.5\) that fails the screening (undefined \(\eta_v^{\text{int}}\)). A population of \(\sim 10^2\)--\(10^3\) autonomous embodied agents with their own energy budgets is required, at the stated power and with preregistration of the protocol (the detection scheme --- § 3.4 main).

\subsection*{S5.4 Summary Lakatosian Passport of the Program}\label{s5.4-summary-lakatosian-passport-of-the-program}

Relative to assembly theory, IIT, teleodynamics, and FEP, the \(\eta_v\)-program is a progressive shift along three independent lines of prediction:

\begin{itemize}
\tightlist
\item
  a monotone correlation of the adaptation rate with \(\eta_v\) (Prediction 1, S5.3);
\item
  a transient pattern of SETI signals as a signature of the externalization-of-payment filter (§ 6 main);
\item
  the existence of a structural threshold \(f^*\) of internalization of an AI system's energy budget (Prediction 3, S5.3).
\end{itemize}

The SETI signature is a progressive prediction \emph{over and above} the nine falsification consequences of § 5 (eight substantively independent lines --- see the Note on numbering), with the status of a speculative extrapolation (§ 6), not a sanctioned operationalization on a calibrated system; it therefore does not count among the consequences of § 5.

The program sustains a progressive regime as long as these predictions admit independent empirical testing.


\nocite{*}
\bibliography{refs}

\end{document}
